

\section{Problemstellung}\label{cha:Problemstellung}
%\section{Problemstatement}\label{cha:problem}

Das Programmierenlernen stellt seit jeher eine erhebliche Hürde in der universitären Ausbildung dar. Wie \textcite[S.~1]{SystematicLiteratureReview2019} in der systematischen Analyse darlegen, weisen Einführungskurse in die Programmierung weltweit  überdurchschnittlich hohe Durchfall- und Abbruchquoten auf. Studierende scheitern oft nicht an der Logik an sich, sondern an der kognitiven Last, die durch die Syntax formaler Sprachen und die Notwendigkeit der abstrakten Problemzerlegung entsteht. Traditionelle Lehrmethoden stoßen hier an ihre Grenzen, da sie oft nicht in der Lage sind, individuelles Feedback in der Frequenz und Tiefe zu liefern, die für das Überwinden dieser initialen Hürden notwendig wäre.

In dieses didaktische Feld treten nun Large Language Models (LLMs) wie ChatGPT, die das Potential haben, die Art und Weise, wie wir Code schreiben und lehren, fundamental zu verändern. Nach \textcite{DennyEra} stehen wir vor einem Paradigmenwechsel. Die Fähigkeit dieser Modelle ausgehend von natürlicher Sprache syntaktisch korrekten Code zu generieren, senkt zudem die Einstiegshürde drastisch \parencite[S.~500-501]{BeckerProgrammingHard}. Doch diese scheinbare Erleichterung birgt ein massives Risiko für die gute wissenschaftliche Praxis und den Kompetenzerwerb: Wenn Studierende die KI nicht als Werkzeug, sondern als Orakel nutzen, entsteht eine „Illusion of Capability"\parencite[S.~116]{PratherRobotsAreHere2023}. Sie produzieren Code, den sie weder verstehen noch warten können, anstatt Problemlösekompetenz zu erwerben.

Ein spezifisches Problem, das in der aktuellen Lage oft übersehen wird, ist die Diskrepanz zwischen der Generierung von isolierten Code-Schnipseln und deren Integration in bestehende, Frameworks. Während LLMs Standardalgorithmen beherrschen, zeigen Studien von \textcite[S.~7]{kazemitabaar_how_2023}, dass Beginner oft daran scheitern, generierten Code korrekt in vorgegebene Strukturen einzubetten oder Kontextfehler wie z.B. halluzinierte Methodenaufrufe zu erkennen.

Die vorliegende Arbeit adressiert genau dieses Feld, indem sie untersucht, wie sich der Einsatz von ChatGPT auf die Bearbeitung einer komplexeren Aufgabe im Rahmen eines vorgegebenen Python Frameworks („Garden-Simulator“) auswirkt. Anstatt sich auf reine Performance-Metriken zu beschränken, kombiniert diese Studie eine qualitative Code-Analyse basierend auf Funktionalitäts- und Stilapekten mit einer Evaluation der User Experience mittels des User Experience Questionnaires (UEQ). Ziel ist es, evidenzbasiert zu klären, ob LLMs in der Hand von Anfängern tatsächlich zu besserer Software führen und in der Bearbeitung der Aufgabe auch bessere Usabibility-Werte erzielen können, als bei einer klassischen Durchführung ohne LLM.
